\chapter{Hyperbolic Three-Space and $\operatorname{PSL}(2,\mathbb C)$}
\label{ch:vii36-h3-psl2c}

\section{The upper half-space model of hyperbolic three-space}

\begin{definition}[Upper half-space]\label{def:vii36-upper}
Hyperbolic three-space is
\[
\mathbb H^3=\{(z,t)\in\mathbb C\times\mathbb R:t>0\}
\]
with metric
\[
ds^2=\frac{|dz|^2+dt^2}{t^2}.
\]
\end{definition}

The metric is conformal to the Euclidean metric on the upper half-space.  Its volume form is
\[
dV=\frac{dx\,dy\,dt}{t^3},\qquad z=x+iy.
\]


\section{Distance, geodesics, and hyperbolic planes}

\begin{theorem}[Distance formula]\label{thm:vii36-distance}
For \(p=(z,t)\) and \(q=(w,s)\),
\[
\cosh d(p,q)=1+\frac{|z-w|^2+(t-s)^2}{2ts}.
\]
\end{theorem}

\begin{theorem}[Geodesics]\label{thm:vii36-geodesics}
The geodesics of \(\mathbb H^3\) are vertical Euclidean lines and Euclidean semicircles orthogonal to the boundary plane \(t=0\).
\end{theorem}

\begin{proposition}[Totally geodesic planes]\label{prop:vii36-planes}
The totally geodesic copies of \(\mathbb H^2\) are vertical Euclidean half-planes and Euclidean hemispheres orthogonal to \(t=0\).
\end{proposition}


\section{Ideal boundary and horospheres}

The natural ideal boundary is
\[
\partial_\infty\mathbb H^3\cong\widehat{\mathbb C}=\mathbb C\cup\{\infty\}.
\]
Each geodesic has two ideal endpoints.  A horosphere based at \(\infty\) is a horizontal Euclidean plane \(t=\text{constant}\); horospheres based at finite boundary points are Euclidean spheres tangent to \(t=0\).

\begin{remark}
The induced metric on a horizontal horosphere \(t=t_0\) is Euclidean up to the constant scale factor \(t_0^{-2}\).
\end{remark}


\section{The hyperboloid model in signature \((3,1)\)}

Let \(\mathbb R^{3,1}\) carry the Lorentz form
\[
\langle x,y\rangle_{3,1}=x_1y_1+x_2y_2+x_3y_3-x_4y_4.
\]
The upper sheet
\[
\mathcal H^3=\{x:\langle x,x\rangle_{3,1}=-1,\ x_4>0\}
\]
with the induced positive-definite metric is another model of hyperbolic three-space.

\begin{theorem}[Lorentzian distance]\label{thm:vii36-lorentz-distance}
For \(x,y\in\mathcal H^3\),
\[
\cosh d(x,y)=-\langle x,y\rangle_{3,1}.
\]
\end{theorem}


\section{Möbius transformations on the Riemann sphere}

Every matrix
\[
A=\begin{pmatrix}a&b\\ c&d\end{pmatrix}\in\operatorname{SL}(2,\mathbb C)
\]
acts on \(\widehat{\mathbb C}\) by
\[
z\longmapsto\frac{az+b}{cz+d}.
\]
The matrices \(A\) and \(-A\) induce the same map, giving \(\operatorname{PSL}(2,\mathbb C)\).

\begin{theorem}[Extension theorem]\label{thm:vii36-extension}
Every Möbius transformation of \(\widehat{\mathbb C}\) extends uniquely to an orientation-preserving isometry of \(\mathbb H^3\), and every orientation-preserving isometry arises this way.  Hence
\[
\operatorname{Isom}^+(\mathbb H^3)\cong\operatorname{PSL}(2,\mathbb C).
\]
\end{theorem}


\section{Classification of individual isometries}

The complex setting adds genuine screw motions.

\begin{definition}[Loxodromic element]\label{def:vii36-loxodromic}
An element of \(\operatorname{PSL}(2,\mathbb C)\) is loxodromic if it has two distinct fixed points on \(\widehat{\mathbb C}\) and, after conjugation, acts as
\[
z\mapsto \lambda z,\qquad |\lambda|\ne1.
\]
\end{definition}

A loxodromic isometry preserves the geodesic joining its two ideal fixed points.  Along that axis it translates and may also rotate.

\begin{definition}[Complex length]\label{def:vii36-complex-length}
For a loxodromic element, the complex length is
\[
L=\ell+i\theta
\]
modulo the usual \(2\pi i\) ambiguity, where \(\ell>0\) is translation length and \(\theta\) is rotation angle about the axis.
\end{definition}


\section{Eigenvalues and trace}

For a lift \(A\in\operatorname{SL}(2,\mathbb C)\) of a loxodromic element, choose eigenvalues
\[
e^{L/2},\qquad e^{-L/2}.
\]
Then
\[
\operatorname{tr}A=2\cosh\frac L2.
\]
This complexifies the trace--translation-length relation from VII/34.

\begin{remark}
In \(\operatorname{PSL}(2,\mathbb C)\), classification by the real number \(|\operatorname{tr}A|\) is no longer sufficient.  The complex trace, or better \((\operatorname{tr}A)^2\), carries rotational information.
\end{remark}


\section{Real Fuchsian groups inside the complex group}

The inclusion
\[
\operatorname{PSL}(2,\mathbb R)\hookrightarrow\operatorname{PSL}(2,\mathbb C)
\]
preserves the boundary circle \(\mathbb R\cup\{\infty\}\).  Its extension to \(\mathbb H^3\) preserves the totally geodesic vertical plane with that ideal boundary.  Thus every Fuchsian group is a special Kleinian group preserving a hyperbolic plane.


\section{Bridge to Kleinian groups}

VII/37 imposes discreteness on subgroups of \(\operatorname{PSL}(2,\mathbb C)\).  The resulting Kleinian groups act simultaneously on \(\mathbb H^3\) and on its Riemann-sphere boundary.  The split between stable boundary action and chaotic accumulation is encoded by the domain of discontinuity and the limit set.

\section{Exercises}

\begin{exercise}\label{exr:vii36-01}
Define the upper half-space model of \(\mathbb H^3\).
\end{exercise}

\begin{hint}
Give the set and metric.
\end{hint}

\begin{solution}
\(\mathbb H^3=\{(z,t)\in\mathbb C\times\mathbb R:t>0\}\) with \(ds^2=(|dz|^2+dt^2)/t^2\).
\end{solution}

\begin{exercise}\label{exr:vii36-02}
What is its hyperbolic volume form?
\end{exercise}

\begin{hint}
Take the determinant of the conformal metric.
\end{hint}

\begin{solution}
\(dV=dx\,dy\,dt/t^3\).
\end{solution}

\begin{exercise}\label{exr:vii36-03}
State the distance formula in upper half-space.
\end{exercise}

\begin{hint}
Use the Euclidean separation divided by heights.
\end{hint}

\begin{solution}
\(\cosh d(p,q)=1+(|z-w|^2+(t-s)^2)/(2ts)\).
\end{solution}

\begin{exercise}\label{exr:vii36-04}
Describe all geodesics.
\end{exercise}

\begin{hint}
Use Euclidean objects orthogonal to the boundary.
\end{hint}

\begin{solution}
They are vertical lines and semicircles orthogonal to \(t=0\).
\end{solution}

\begin{exercise}\label{exr:vii36-05}
Describe totally geodesic planes.
\end{exercise}

\begin{hint}
Raise the dimension of the geodesic description by one.
\end{hint}

\begin{solution}
They are vertical half-planes and hemispheres orthogonal to the boundary plane.
\end{solution}

\begin{exercise}\label{exr:vii36-06}
Identify the ideal boundary of \(\mathbb H^3\).
\end{exercise}

\begin{hint}
Compactify the boundary plane by infinity.
\end{hint}

\begin{solution}
\(\partial_\infty\mathbb H^3\cong\widehat{\mathbb C}\).
\end{solution}

\begin{exercise}\label{exr:vii36-07}
What is a horosphere based at infinity?
\end{exercise}

\begin{hint}
Use a level set of height.
\end{hint}

\begin{solution}
A horizontal plane \(t=t_0\).
\end{solution}

\begin{exercise}\label{exr:vii36-08}
What is a horosphere based at a finite boundary point?
\end{exercise}

\begin{hint}
Think of Euclidean spheres tangent to the boundary.
\end{hint}

\begin{solution}
A Euclidean sphere tangent to the plane \(t=0\) at that boundary point.
\end{solution}

\begin{exercise}\label{exr:vii36-09}
Write the Lorentz form used in the hyperboloid model.
\end{exercise}

\begin{hint}
Use three plus signs and one minus sign.
\end{hint}

\begin{solution}
\(\langle x,y\rangle_{3,1}=x_1y_1+x_2y_2+x_3y_3-x_4y_4\).
\end{solution}

\begin{exercise}\label{exr:vii36-10}
Define the hyperboloid model.
\end{exercise}

\begin{hint}
Take the upper unit timelike sheet.
\end{hint}

\begin{solution}
\(\mathcal H^3=\{x:\langle x,x\rangle_{3,1}=-1,\ x_4>0\}\).
\end{solution}

\begin{exercise}\label{exr:vii36-11}
State the hyperboloid distance formula.
\end{exercise}

\begin{hint}
Use the Lorentz inner product.
\end{hint}

\begin{solution}
\(\cosh d(x,y)=-\langle x,y\rangle_{3,1}\).
\end{solution}

\begin{exercise}\label{exr:vii36-12}
How does \(A\in\operatorname{SL}(2,\mathbb C)\) act on the Riemann sphere?
\end{exercise}

\begin{hint}
Use fractional linear transformation.
\end{hint}

\begin{solution}
\(z\mapsto(az+b)/(cz+d)\).
\end{solution}

\begin{exercise}\label{exr:vii36-13}
Why do \(A\) and \(-A\) define the same Möbius transformation?
\end{exercise}

\begin{hint}
Both numerator and denominator change sign.
\end{hint}

\begin{solution}
The common sign cancels in the quotient, so the action factors through PSL.
\end{solution}

\begin{exercise}\label{exr:vii36-14}
Identify \(\operatorname{Isom}^+(\mathbb H^3)\).
\end{exercise}

\begin{hint}
Use the extension theorem.
\end{hint}

\begin{solution}
\(\operatorname{Isom}^+(\mathbb H^3)\cong\operatorname{PSL}(2,\mathbb C)\).
\end{solution}

\begin{exercise}\label{exr:vii36-15}
How many ideal fixed points does a loxodromic element have?
\end{exercise}

\begin{hint}
Use its normal form.
\end{hint}

\begin{solution}
Exactly two distinct fixed points.
\end{solution}

\begin{exercise}\label{exr:vii36-16}
What geometric object joins those fixed points?
\end{exercise}

\begin{hint}
There is a unique geodesic with given ideal endpoints.
\end{hint}

\begin{solution}
The invariant geodesic axis.
\end{solution}

\begin{exercise}\label{exr:vii36-17}
What two motions can a loxodromic isometry perform along its axis?
\end{exercise}

\begin{hint}
Complex length has real and imaginary parts.
\end{hint}

\begin{solution}
A translation along the axis and a rotation about it.
\end{solution}

\begin{exercise}\label{exr:vii36-18}
Define complex length.
\end{exercise}

\begin{hint}
Separate translation and rotation.
\end{hint}

\begin{solution}
\(L=\ell+i\theta\), with \(\ell>0\) translation length and \(\theta\) rotation angle modulo \(2\pi\).
\end{solution}

\begin{exercise}\label{exr:vii36-19}
Give the eigenvalues associated with complex length \(L\).
\end{exercise}

\begin{hint}
Use determinant one.
\end{hint}

\begin{solution}
They may be chosen as \(e^{L/2}\) and \(e^{-L/2}\).
\end{solution}

\begin{exercise}\label{exr:vii36-20}
Write the complex trace-length relation.
\end{exercise}

\begin{hint}
Add the two eigenvalues.
\end{hint}

\begin{solution}
\(\operatorname{tr}A=2\cosh(L/2)\).
\end{solution}

\begin{exercise}\label{exr:vii36-21}
Why is \(|\operatorname{tr}A|\) alone insufficient in \(\operatorname{PSL}(2,\mathbb C)\)?
\end{exercise}

\begin{hint}
Rotational data are complex.
\end{hint}

\begin{solution}
Different complex traces can have the same modulus while encoding different rotation angles.
\end{solution}

\begin{exercise}\label{exr:vii36-22}
What subset of the ideal boundary is preserved by \(\operatorname{PSL}(2,\mathbb R)\)?
\end{exercise}

\begin{hint}
Real coefficients preserve the real projective line.
\end{hint}

\begin{solution}
The circle \(\mathbb R\cup\{\infty\}\).
\end{solution}

\begin{exercise}\label{exr:vii36-23}
What subset of \(\mathbb H^3\) is then preserved?
\end{exercise}

\begin{hint}
Take the totally geodesic plane with that ideal boundary.
\end{hint}

\begin{solution}
The vertical totally geodesic copy of \(\mathbb H^2\).
\end{solution}

\begin{exercise}\label{exr:vii36-24}
What is the next step from Fuchsian to Kleinian groups?
\end{exercise}

\begin{hint}
Change the ambient group but keep discreteness.
\end{hint}

\begin{solution}
Study discrete subgroups of \(\operatorname{PSL}(2,\mathbb C)\) and their action on both \(\mathbb H^3\) and \(\widehat{\mathbb C}\).
\end{solution}

\section{Solved problem dossiers}

\begin{problem}[Vertical distance]\label{prob:vii36-01}
Compute the distance between \((0,t_1)\) and \((0,t_2)\) on a vertical geodesic.
\end{problem}

\begin{solution}
Restricting the metric gives \(ds=dt/t\).  Hence
\[
d((0,t_1),(0,t_2))=\left|\int_{t_1}^{t_2}\frac{dt}{t}\right|=\left|\log\frac{t_2}{t_1}\right|.
\]
\end{solution}

\begin{problem}[Checking the distance formula vertically]\label{prob:vii36-02}
Verify the general distance formula when \(z=w\).
\end{problem}

\begin{solution}
The formula gives
\[
\cosh d=1+\frac{(t-s)^2}{2ts}=\frac{t^2+s^2}{2ts}=\cosh\left(\log\frac ts\right),
\]
consistent with the vertical-distance computation.
\end{solution}

\begin{problem}[A totally geodesic plane]\label{prob:vii36-03}
Show that the subset \(\{(x,t):x\in\mathbb R,t>0\}\subset\mathbb H^3\) is isometric to the usual upper half-plane.
\end{problem}

\begin{solution}
Set the second horizontal coordinate equal to zero.  The induced metric is \((dx^2+dt^2)/t^2\), exactly the Poincaré upper-half-plane metric.  It is therefore a totally geodesic copy of \(\mathbb H^2\).
\end{solution}

\begin{problem}[Dilation isometry]\label{prob:vii36-04}
Show that \((z,t)\mapsto(az,at)\) for \(a>0\) is an isometry of upper half-space.
\end{problem}

\begin{solution}
The numerator \(|dz|^2+dt^2\) scales by \(a^2\), while the denominator \(t^2\) also scales by \(a^2\).  Their ratio is unchanged.
\end{solution}

\begin{problem}[Boundary action of a dilation]\label{prob:vii36-05}
Identify the boundary Möbius map induced by the previous dilation and its fixed points.
\end{problem}

\begin{solution}
On \(\widehat{\mathbb C}\) it is \(z\mapsto az\).  Its fixed points are \(0\) and \(\infty\), the ideal endpoints of the vertical axis.
\end{solution}

\begin{problem}[A pure screw motion]\label{prob:vii36-06}
Consider \(z\mapsto re^{i\theta}z\) with \(r>1\).  Describe its extension geometrically.
\end{problem}

\begin{solution}
It preserves the vertical axis joining \(0\) to \(\infty\).  It translates by \(\log r\) along that axis and rotates through angle \(\theta\) about it, so its complex length is \(\log r+i\theta\) up to convention.
\end{solution}

\begin{problem}[Trace from complex length]\label{prob:vii36-07}
Suppose a lift has eigenvalues \(e^{L/2}\) and \(e^{-L/2}\).  Derive the trace relation.
\end{problem}

\begin{solution}
Adding the eigenvalues gives
\[
\operatorname{tr}A=e^{L/2}+e^{-L/2}=2\cosh(L/2).
\]
\end{solution}

\begin{problem}[Why PSL rather than SL]\label{prob:vii36-08}
Explain geometrically why the isometry group is PSL rather than SL.
\end{problem}

\begin{solution}
The central matrices \(I\) and \(-I\) induce the same fractional linear transformation on the boundary.  By uniqueness of extension, they induce the same interior isometry.  Thus the kernel is \(\{\pm I\}\).
\end{solution}

\begin{problem}[Horosphere metric]\label{prob:vii36-09}
Compute the induced metric on the horosphere \(t=t_0\).
\end{problem}

\begin{solution}
Setting \(dt=0\) gives \(ds^2=|dz|^2/t_0^2\).  This is a Euclidean metric scaled by the constant factor \(t_0^{-2}\).
\end{solution}

\begin{problem}[Boundary circle and hyperbolic plane]\label{prob:vii36-10}
Explain why a generalized circle in \(\widehat{\mathbb C}\) determines a unique totally geodesic plane in \(\mathbb H^3\).
\end{problem}

\begin{solution}
If the circle passes through \(\infty\), it is a Euclidean line and bounds a vertical half-plane.  Otherwise it is a Euclidean circle and bounds the hemisphere orthogonal to \(t=0\).  These are exactly the totally geodesic planes.
\end{solution}

\begin{problem}[Fuchsian subgroup in three dimensions]\label{prob:vii36-11}
Let \(\Gamma<\operatorname{PSL}(2,\mathbb R)\).  Show that its three-dimensional action preserves a copy of \(\mathbb H^2\).
\end{problem}

\begin{solution}
Its boundary action preserves \(\mathbb R\cup\{\infty\}\).  The unique totally geodesic plane with this ideal boundary is the vertical plane over the real axis, hence it is \(\Gamma\)-invariant.
\end{solution}

\begin{problem}[Two endpoints determine an axis]\label{prob:vii36-12}
Show that two distinct ideal points determine a unique geodesic in upper half-space.
\end{problem}

\begin{solution}
If one endpoint is infinity, the geodesic is the vertical line over the finite endpoint.  If both are finite, there is a unique Euclidean circle through them with center on the boundary plane and orthogonal to it; its upper semicircle is the required geodesic.
\end{solution}

\section{Challenge problems}

\begin{problem}[Complex trace ambiguity]\label{prob:vii36-ch01}
Show that \((\operatorname{tr}A)^2\) is well defined for an element of \(\operatorname{PSL}(2,\mathbb C)\), although \(\operatorname{tr}A\) itself is only defined up to sign.
\end{problem}

\begin{solution}
The two lifts are \(A\) and \(-A\).  Their traces differ by a sign, so squaring removes the ambiguity.  Thus \((\operatorname{tr}A)^2\) is an invariant of the PSL element.
\end{solution}

\begin{problem}[Loxodromic dynamics on the boundary]\label{prob:vii36-ch02}
For \(g(z)=\lambda z\) with \(|\lambda|>1\), determine the forward and backward limits of every \(z\ne0,\infty\).
\end{problem}

\begin{solution}
As \(n\to+\infty\), \(|\lambda^n z|\to\infty\), so \(g^nz\to\infty\).  As \(n\to+\infty\), \(g^{-n}z=\lambda^{-n}z\to0\).  Thus \(\infty\) is attracting and \(0\) repelling.
\end{solution}

\begin{problem}[Recovering the two-dimensional theory]\label{prob:vii36-ch03}
Explain how VII/34 is recovered by restricting to real matrices with real eigenvalues.
\end{problem}

\begin{solution}
For real hyperbolic matrices the rotation angle is \(0\), so complex length reduces to the real translation length \(L=\ell\).  The trace relation becomes \(|\operatorname{tr}A|=2\cosh(\ell/2)\), exactly the two-dimensional formula.
\end{solution}

\begin{problem}[Boundary data determine the interior isometry]\label{prob:vii36-ch04}
Why is uniqueness in the Möbius extension theorem plausible from hyperbolic geometry?
\end{problem}

\begin{solution}
An orientation-preserving isometry maps geodesics to geodesics, and each geodesic is determined by its two ideal endpoints.  If two isometries have the same boundary action, they agree on every geodesic endpoint pair and hence on the geodesic configuration; this forces agreement throughout \(\mathbb H^3\).
\end{solution}

\begin{problem}[From one isometry to a discrete group]\label{prob:vii36-ch05}
What new phenomenon must be controlled when replacing a single loxodromic isometry by a subgroup of \(\operatorname{PSL}(2,\mathbb C)\)?
\end{problem}

\begin{solution}
One must control accumulation of infinitely many group elements and orbits.  Discreteness gives a good action in the interior, but boundary orbits can accumulate on complicated closed sets.  This leads to the Kleinian limit set and domain of discontinuity.
\end{solution}
